% Source-aligned English for printed p.204.
\clearpage
\SourcePage{204}

\LayerTarget{tohoku:remark:condition-D-basic-properties:p1}
\noindent Zariski topology and $G$ is a finite group of automorphisms of $X$,
condition (D) does \textit{not} hold (unless $G$ acts trivially).

\LayerTarget{tohoku:theorem:5.3.1}
\textsc{Theorem 5.3.1.} \textit{Suppose that condition (D) holds.
Let $A$ be an abelian $G$-sheaf on $X$, let $y\in Y$, and let $x$ be an
element of $f^{-1}(y)$. There are then canonical isomorphisms}

\LayerTarget{tohoku:equation:5.3.1}
\[
\mathbf H^n(G,A)(y)=H^n(G_x,A(x)).
\tag{5.3.1}
\]

\LayerTarget{tohoku:proof:theorem-5.3.1}
Taking $V_x$ as in the statement of condition (D), and moreover taking $V_x$
open and stable under $G_x$, we may apply the corollary to Proposition 5.2.1. This
reduces us to the case $G=G_x$, $X=V_x$, and hence to $G$ \textit{finite}.
Let us show that $R^qf_*(A)=0$ for $q>0$. For every $y\in Y$ we have

\LayerTarget{tohoku:equation:5.3.1-stalk-limit}
\[
R^qf_*(A)(y)=\mathop{\underset{\to}{\mathrm{lim.}}} H^q(f^{-1}(U),A),
\]
where the limit is taken over the open neighborhoods $U$ of $y$. With $x$ and
$V_x$ chosen as above, and restricting to $U\subset f(V_x)$, we see that
$f^{-1}(U)$ is a finite union of pairwise disjoint open sets $g\mathbin{.}U_x$
($g\in G/G_x$), where $U_x=V_x\cap f^{-1}(U)$. Thus
$H^q(f^{-1}(U),A)$ is a direct sum of finitely many groups isomorphic to
$H^q(U_x,A)$. As $U$ ranges over a fundamental system of neighborhoods of
$y$, $U_x$ ranges over a fundamental system of neighborhoods of $x$; hence
for $R^qf_*(A)(y)$ we obtain a direct sum of finitely many groups isomorphic to
$\mathop{\underset{\to}{\lim}} H^q(U_x,A)$. The latter is zero for $q>0$ by Lemma 3.8.2.

Thus indeed $R^qf_*(A)=0$ for $q>0$. By Proposition 5.2.2, the left-hand side
of (5.3.1) is therefore identified with $\mathbf H^n(G,f_*(A))(y)$, which,
since $G$ is finite, is equal to $H^n(G,f_*(A)(y))$ (as noted in 4.4).
Recalling that we could suppose $X=V_x$, so that $f^{-1}(y)$ is reduced to
$x$, we have $f_*(A)(y)=A(x)$, which completes the proof.

\LayerTarget{tohoku:corollary:5.3.1:1}
\textsc{Corollary.} \textit{Suppose that for every $x\in X$ the
order $n_x$ of the stabilizer $G_x$ is such that multiplication by $n_x$ is
an automorphism of the group $A(x)$. Then $H^*(X;G,A)=H^*(Y,A^G)$;
consequently $H^*(Y,A^G)$ is the abutment of a spectral sequence whose
initial term is $II_2^{p,q}=H^p(G,H^q(X,A))$.}

\LayerTarget{tohoku:proof:corollary-5.3.1:1}
Indeed, Theorem 5.3.1 gives $\mathbf H^q(G,A)=0$ for $q>0$, so the hypotheses
of Proposition 5.2.3 apply. The preceding corollary is especially useful when
$A$ is a sheaf of vector spaces over a field of characteristic $p$ that
divides none of the $n_x$ (as is always the case in characteristic $0$). When
there is no ground field, one may use the following variant.

\LayerTarget{tohoku:corollary:5.3.1:2}
\textsc{Corollary 2.} \textit{Suppose that the orders $n_x$ of the
$G_x$ admit a least common multiple $n$; let $P$ be the set of prime divisors
of $n$, and let ${\mathbf C}(P)$ be the category of abelian groups each of whose elements
has finite order with no prime divisors other than members $p\in P$. Then, in
the terminology introduced in \textup{1.11}, $H^*(X;G,A)$ is isomorphic modulo ${\mathbf C}(P)$
to $H^*(Y,A^G)$; consequently the latter is, modulo ${\mathbf C}(P)$, the abutment of a
cohomological spectral sequence whose initial term is $H^p(G,H^q(X,A))$.}

\LayerTarget{tohoku:proof:corollary-5.3.1:2:p1}
Indeed, reduce modulo ${\mathbf C}(P)$ the spectral sequences of Theorem 5.2.1.

\TohokuContinuation{205}
